% 第七章 系统测试与效果分析

\chapter{系统测试与效果分析}\label{ch:system_testing}


\section{验证范围与实验环境}

FreeFlow 同时涉及桌面端、Web2.5 服务器和智能合约三部分。桌面端交互、IPFS 网关和钱包签名已在第\ref{ch:core_modules}章中详细阐述，本章节将可稳定复现的验证工作集中在两部分：一是智能合约功能正确性的自动化测试，二是链上资源检索排序模块的离线评估。

实验所使用的硬件、软件与数据规模配置如表\ref{tab:test_environment}所示。

\begin{table}[htbp]
    \linespread{1.5}
    \zihao{5}
    \songti
    \centering
    \caption{实验环境配置}\label{tab:test_environment}
    \begin{tabular}{|l|p{10cm}|}
        \hline
        类别      & 环境说明                           \\ \hline
        硬件设备    & Apple M4，16 GB 内存              \\ \hline
        操作系统    & macOS 26.1（arm64）              \\ \hline
        后端运行时   & Node.js v20.18.0，pnpm 10.0.0   \\ \hline
        桌面端技术栈  & Electron、React、TypeScript      \\ \hline
        后端技术栈   & Node.js、Express、Prisma         \\ \hline
        合约测试工具链 & Hardhat 2.28.6，Solidity 0.8.26 \\ \hline
        排序评估集规模 & 6000条候选资源，162个查询               \\ \hline
        参考时间点   & 2026-04-22 12:00:00 UTC        \\ \hline
    \end{tabular}
\end{table}


\section{智能合约自动化测试结果}

智能合约测试基于Hardhat测试框架完成，围绕作品发布、购买授权、访问控制和销售配置修改等核心场景设计自动化用例。测试总计4项，全部通过，单次执行耗时不到1 s。测试结果如表\ref{tab:contract_tests}所示。

\begin{table}[htbp]
    \linespread{1.45}
    \zihao{5}
    \songti
    \centering
    \caption{智能合约自动化测试结果}\label{tab:contract_tests}
    \begin{tabular}{|l|p{8.3cm}|l|}
        \hline
        编号  & 测试内容                            & 结果 \\ \hline
        C01 & 付费作品发布后，用户支付正确金额可获得ERC-1155访问凭证 & 通过 \\ \hline
        C02 & 公开作品发布后，未购买用户默认拥有访问权限           & 通过 \\ \hline
        C03 & ERC-1155访问凭证不可在普通用户之间转让         & 通过 \\ \hline
        C04 & 只有创作者本人可以修改作品销售配置               & 通过 \\ \hline
    \end{tabular}
\end{table}

该结果表明，\texttt{PlatformHub}与\texttt{MusicAccess1155}的协同配合能正确运行；同时，创作者侧收益分配能够按预设流程导入\texttt{RoyaltySplitter}，但是生产环境的安全测试和大规模链上压力测试仍需在后续工作中继续完成。


\section{检索排序实验设置}

离线检索评估集是围绕第\ref{ch:search_ranking}章所提出的链上资源排序目标来进行构建的。该评估集基于公开 MusicBrainz sample 2026-04-01 数据库生成，先选取 48 个发行簇作为核心相关样本，再从整库中抽取干扰项补齐候选集，最终形成了 6000 条候选资源，其中核心相关资源为 815 条、干扰项为 5185 条。查询集共 162 条，用于覆盖精确检索、模糊检索、链上身份定位以及输入扰动等场景。对于标题、专辑、艺人、别名和多词组合等多结果查询，采用 3/2/1 相关等级标注；对于 Token ID、合约地址、CID、\texttt{resourceKey}和交易哈希等高确定性定位查询，则采用浅层 spot-check 标注。

\begin{table}[htbp]
    \linespread{1.45}
    \zihao{5}
    \songti
    \centering
    \caption{检索评估查询类型分布}\label{tab:ranking_query_categories}
    \begin{tabular}{|l|c|l|}
        \hline
        查询类型           & 数量 & 说明                                      \\ \hline
        标题精确匹配         & 24 & 直接定位单首作品                                \\ \hline
        专辑查询           & 24 & 验证同专辑资源排序                               \\ \hline
        艺术家查询          & 18 & 返回同创作者多首作品                              \\ \hline
        标题+艺人组合查询      & 18 & 验证标题与艺人联合意图                            \\ \hline
        多词组合查询         & 18 & 验证艺人与专辑联合检索                            \\ \hline
        发行别名查询         & 10 & 验证 release alias 检索能力                   \\ \hline
        艺人别名查询         & 10 & 验证 artist alias 检索能力                    \\ \hline
        拼写扰动查询         & 12 & 如\texttt{Introductiox}                       \\ \hline
        紧凑输入/空格扰动      & 10 & 如\texttt{MorningDips}                      \\ \hline
        全角输入           & 4  & 验证 NFKC 归一化                            \\ \hline
        Token ID 查询    & 3  & 验证链上标识精确定位                              \\ \hline
        合约地址查询         & 3  & 验证音乐合约地址检索                              \\ \hline
        CID 查询         & 3  & 验证 IPFS 内容寻址检索                         \\ \hline
        resourceKey 查询 & 3  & 验证链上资源键精确定位                             \\ \hline
        交易哈希查询         & 2  & 验证发布交易定位                                \\ \hline
    \end{tabular}
\end{table}

实验内部共设置 4 组方法，其中“文本+链上身份字段”属于中间消融结果，主要用于内部调参与分析。为了让正文展示更加紧凑，本文在表格与图中仅保留 3 组代表性方法：数据库字段包含匹配基线、仅文本评分方法，以及完整的\texttt{chain\_resource\_rank\_v1}排序算法。评价指标方面，囊括了 Precision@10、MRR、nDCG@10、零结果率以及平均响应时间。


\section{检索效果对比实验}

表\ref{tab:ranking_metrics_actual}和图\ref{fig:ranking_metrics_bar}给出了 3 种代表性方法在 162 个查询上的总体表现。由于评估集同时包含较多单目标定位查询和多结果查询，Precision@10 的绝对值更容易受到相关结果数量上限的影响，因此本节重点结合 MRR 和 nDCG@10 分析排序质量。

\begin{table}[htbp]
    \linespread{1.45}
    \zihao{5}
    \songti
    \centering
    \caption{检索排序指标对比结果}\label{tab:ranking_metrics_actual}
    \begin{tabular}{|p{3.2cm}|c|c|c|c|c|}
        \hline
        方法          & Precision@10 & MRR    & nDCG@10 & 零结果率    & 平均响应时间/ms \\ \hline
        数据库字段包含匹配基线 & 0.1377       & 0.4681 & 0.4334  & 47.53\% & 23.2932   \\ \hline
        仅文本评分       & 0.2105       & 0.6992 & 0.6014  & 10.49\% & 278.6325  \\ \hline
        完整排序算法      & 0.1963       & 0.7253 & 0.6480  & 3.70\%  & 297.7463  \\ \hline
    \end{tabular}
\end{table}

\begin{figure}[htbp]
    \centering
    \begin{tikzpicture}
        \begin{axis}[
            ybar,
            width=0.98\textwidth,
            height=8.1cm,
            bar width=9pt,
            ymin=0,
            ymax=105,
            ylabel={指标值（\%）},
            symbolic x coords={sql_contains,text_only,full_ranker},
            xtick=data,
            xticklabels={字段包含匹配基线,仅文本评分,完整排序算法},
            xticklabel style={rotate=18,anchor=east,font=\zihao{6}},
            ymajorgrids=true,
            legend style={at={(0.5,1.05)},anchor=south,legend columns=3,font=\zihao{6}},
            enlarge x limits=0.08,
        ]
            \addplot coordinates {(sql_contains,13.7654) (text_only,21.0494) (full_ranker,19.6296)};
            \addplot coordinates {(sql_contains,46.8056) (text_only,69.9177) (full_ranker,72.5279)};
            \addplot coordinates {(sql_contains,43.3390) (text_only,60.1417) (full_ranker,64.7998)};
            \legend{Precision@10,MRR,nDCG@10}
        \end{axis}
    \end{tikzpicture}
    \caption{不同排序策略的检索质量对比}
    \label{fig:ranking_metrics_bar}
\end{figure}

从实验结果可以看出，单纯依赖数据库字段包含匹配虽然开销最低，但在拼写扰动、别名检索和链上复杂身份查询场景中会出现明显漏召回，零结果率达到 47.53\%。仅文本评分方法已经能够处理一部分模糊查询与紧凑输入，Precision@10 达到 0.2105，但由于缺少链上身份信号和辅助排序能力，在首位命中质量上仍存在不足，其 MRR 和 nDCG@10 分别只有 0.6992 和 0.6014。

完整排序算法在综合文本相关性、链上身份字段、新鲜度和热度信号后，把零结果率进一步压低到 3.70\%，同时把 MRR 提升到 0.7253、nDCG@10 提升到 0.6480，说明目标资源更容易稳定地出现在结果前列。虽然它在 Precision@10 上略低于仅文本评分方法，但在首位命中和分级排序质量上表现更好，这与资源发现模块更强调“前几位是否可靠命中目标”的使用目标是一致的。


\section{典型查询案例分析}

完整排序算法除了输出结果列表外，还会返回排序原因，便于分析命中特征与排序依据。表\ref{tab:ranking_case_analysis}给出了3类典型查询的实验结果。

\begin{table}[htbp]
    \linespread{1.45}
    \zihao{5}
    \songti
    \centering
    \caption{典型查询案例分析}\label{tab:ranking_case_analysis}
    \begin{tabular}{|p{3.6cm}|p{3.6cm}|p{5.2cm}|}
        \hline
        查询输入                               & 首位结果                      & 主要排序原因                          \\ \hline
        \texttt{Introductiox}              & Introduction            & 文本相似度约92\%，能够纠正尾部单字符错误               \\ \hline
        \texttt{MorningDips}               & Morning Dips            & 标题经紧凑化后精确匹配，能够处理去空格输入             \\ \hline
        resourceKey精确查询                    & Houston（日文曲目）          & resourceKey精确匹配，同时命中音乐合约地址             \\ \hline
    \end{tabular}
\end{table}

上述结果表明，排序模块不仅能够处理传统文本检索，还能处理链上身份字段定位以及带空格、大小写和轻微拼写误差的输入。排序原因字段使系统具备较强的可解释性，也便于后续继续调权和排查误召回问题。


\section{性能测试结果}

为了观察完整排序算法在不同候选规模下的开销变化，本文分别在 1000、5000 和 6000 条候选资源规模下统计平均响应时间与 95\% 分位响应时间。结果如表\ref{tab:ranking_latency}和图\ref{fig:ranking_latency_curve}所示。

\begin{table}[htbp]
    \linespread{1.45}
    \zihao{5}
    \songti
    \centering
    \caption{不同候选规模下的排序耗时}\label{tab:ranking_latency}
    \begin{tabular}{|c|c|c|}
        \hline
        候选数量 & 平均响应时间/ms & 95\%分位响应时间/ms \\ \hline
        1000 & 46.00     & 83.93         \\ \hline
        5000 & 203.55    & 367.19        \\ \hline
        6000 & 252.50    & 449.40        \\ \hline
    \end{tabular}
\end{table}

\begin{figure}[htbp]
    \centering
    \begin{tikzpicture}
        \begin{axis}[
            width=0.94\textwidth,
            height=7.8cm,
            xlabel={候选数量},
            ylabel={响应时间/ms},
            xmin=900,
            xmax=6100,
            ymin=0,
            ymax=500,
            ymajorgrids=true,
            legend style={at={(0.5,1.05)},anchor=south,legend columns=2,font=\zihao{6}},
            tick label style={font=\zihao{6}},
            label style={font=\zihao{6}},
        ]
            \addplot+[mark=*,line width=1pt] table[x=candidate_count,y=avg_latency_ms,col sep=comma]{figures/data/search_latency_candidates.csv};
            \addplot+[mark=square*,line width=1pt] table[x=candidate_count,y=p95_latency_ms,col sep=comma]{figures/data/search_latency_candidates.csv};
            \legend{平均响应时间,95\%分位响应时间}
        \end{axis}
    \end{tikzpicture}
    \caption{完整排序算法在不同候选规模下的延迟变化}
    \label{fig:ranking_latency_curve}
\end{figure}

结果显示，随着候选数量增加，平均响应时间近似线性增长，这与第\ref{ch:search_ranking}章中的复杂度分析一致。在 6000 条候选资源下，完整排序算法的平均响应时间为 252.50 ms，95\% 分位响应时间为 449.40 ms。在线场景中，系统仍会先通过数据库字段召回与候选集截断来限制参与排序的规模，因此该延迟结果主要用于说明算法在更大候选集合上的增长趋势。


\section{结果讨论与有效性说明}

综合本章结果，可以得出两点结论。第一，系统当前已经在可复现范围内验证了链上发布、购买和访问控制的关键业务链路，说明合约层设计足以支持本文设计的凭证化和访问授权。第二，检索排序模块在冷启动场景下具备较好的工程可用性；在基于公开曲库构建的离线基准上，完整排序算法相较于基础字段匹配和仅文本评分方法，能够更有效的召回相关项。

同时，本章也有明确边界。桌面端的登录、上传、详情展示、播放、评论和下载等功能已经完成实现与联调，但检索方法、生产级链上场景尚未形成大规模自动化测试。FreeFlow 完成了基于区块链，面向音乐场景的混合架构实现验证，生产环境中的完整平台成熟度仍有待进一步提升。
